\documentclass[11pt,a4paper]{extbook}

\usepackage{fontspec}
\usepackage{polyglossia}
\setdefaultlanguage{russian}
\setotherlanguage{english}

\usepackage{geometry}
\usepackage{graphicx}
\usepackage{hyperref}
\usepackage{microtype}
\usepackage{amsmath}
\usepackage{amssymb}
\usepackage{mathtools}
\usepackage{booktabs}
\usepackage{longtable}
\usepackage{array}
\usepackage{enumitem}
\usepackage{fancyhdr}
\usepackage{xcolor}
\usepackage{titlesec}
\usepackage[most]{tcolorbox}

\defaultfontfeatures{Ligatures=TeX}
\setmainfont{DejaVu Serif}
\setsansfont{DejaVu Sans}
\setmonofont{DejaVu Sans Mono}[Scale=0.88]
\DeclareMathSizes{11}{13}{10}{8}

\geometry{left=2.4cm,right=2.4cm,top=2.5cm,bottom=2.5cm}
\setlength{\parindent}{0pt}
\setlength{\parskip}{0.45em}
\setlength{\jot}{5pt}
\tolerance=2000
\emergencystretch=2em

\hypersetup{
  colorlinks=true,
  linkcolor=black,
  urlcolor=blue,
  citecolor=black
}

\titleformat{\chapter}[hang]{\Large\bfseries}{\thechapter.}{16pt}{\Large\bfseries}
\titleformat{\section}[hang]{\large\bfseries}{\thesection.}{12pt}{\large\bfseries}
\titleformat{\subsection}[hang]{\normalsize\bfseries}{\thesubsection.}{10pt}{\normalsize\bfseries}

\setlength{\headheight}{22pt}
\pagestyle{fancy}
\fancyhf{}
\lhead{\small\nouppercase{\leftmark}}
\cfoot{\thepage}

\newtcolorbox{takeaway}[1][Вывод]{colback=gray!5,colframe=gray!50,title={#1},fonttitle=\bfseries,breakable}
\newtcolorbox{methodbox}[1][Методический акцент]{colback=blue!4,colframe=blue!40,title={#1},fonttitle=\bfseries,breakable}
\newtcolorbox{riskbox}[1][Риск интерпретации]{colback=red!4,colframe=red!35,title={#1},fonttitle=\bfseries,breakable}
\newtcolorbox{openbox}[1][Открытый вопрос]{colback=orange!8,colframe=orange!60!black,title={#1},fonttitle=\bfseries,breakable}

\newcommand{\score}{\mathrm{Score}}
\newcommand{\PR}{\mathrm{PR}}
\newcommand{\simf}{\mathrm{sim}}
\newcommand{\incoming}{\mathrm{incoming}}
\newcommand{\retained}{\mathrm{retained}}
\newcommand{\degout}{k_{\mathrm{out}}}
\newcommand{\sout}{s_{\mathrm{out}}}
\newcommand{\fanin}{\mathrm{fanin}}
\newcommand{\capacity}{C}

\title{Reports 2--9\\[0.4em]\large Развернутый исследовательский разбор серии экспериментов по графовой памяти RIN}
\author{Рабочая ChatGPT-версия исследовательского свода}
\date{18 августа 2026 г.}

\begin{document}

\frontmatter
\maketitle

\begin{center}
\begin{minipage}{0.88\textwidth}\small
Этот проект собирает в одном месте исследовательскую линию от \texttt{exp2} до
\texttt{exp9a}. Жанр текста -- не набор коротких отчетов, а последовательный
курс-разбор: что именно проверял каждый эксперимент, почему он возник именно в
этот момент, какую альтернативу он пытался отличить от конкурирующей гипотезы,
какой стенд был собран, что показали результаты и как они изменили дальнейший
план работ. Отдельно проговариваются пределы каждого результата: где мы имеем
дело с полноценным эмпирическим опровержением, где с диагностикой, а где лишь с
хорошо поставленной развилкой для следующего шага.
\end{minipage}
\end{center}

\tableofcontents

\mainmatter

\chapter*{Предисловие}
\addcontentsline{toc}{chapter}{Предисловие}

Этот текст нужен не для того, чтобы еще раз кратко пересказать ноутбуки, а для
того, чтобы собрать из них единую исследовательскую историю. В серии
\texttt{exp2--exp9a} происходило не линейное ``улучшение метрики'', а медленное
смещение вопроса. Сначала задача выглядела как инженерная: построить эмбрион
памяти из реальных источников и проверить, что по нему вообще можно пустить
сигнал. Потом стало ясно, что сам факт работоспособности пайплайна почти ничего
не говорит о качестве retrieval. Затем выяснилось, что недостаточно просто
добавить семантический gate: на уровне топологии и структуры сходимости граф
ведет себя иначе, чем хотелось бы. После этого линия распалась на несколько
ветвей:

\begin{enumerate}[leftmargin=2em]
\item ветвь локальной физики сигнала: что удерживает входящий поток и чем
отличается узел-ответ от узла-хаба;
\item ветвь глобальных поправок: как отделить яркость, вызванную конкретным
запросом, от яркости, вызванной общей структурной славой узла;
\item ветвь переноса на реальный эмбрион: сохраняются ли выводы, полученные на
читабельных золотых мирах;
\item ветвь транспортной реконструкции: не в скоринге ли проблема, а в том, что
сигнал вообще идет не туда или не доходит до цели;
\item ветвь качества самой геометрии сходства: не слишком ли грязен косинус в
узком конусе NumberBatch-пространства;
\item ветвь реконструкции недостающих путей: можно ли достраивать граф
семантическим каналом, не превратив его в обычный векторный поиск.
\end{enumerate}

Во всех этих ветвях важно различать три уровня утверждений.

\begin{enumerate}[leftmargin=2em]
\item \textbf{Положительный результат.} Эксперимент убедительно показывает, что
некоторый механизм действительно выполняет свою локальную работу.
\item \textbf{Отрицательный результат.} Эксперимент убедительно показывает, что
некоторая популярная надежда не срабатывает.
\item \textbf{Диагностический результат.} Эксперимент не дает победителя, но
сильно сужает пространство следующих гипотез.
\end{enumerate}

Серия \texttt{exp2--exp9a} особенно богата именно третьим типом результата.
Это не недостаток. Напротив, в архитектурной работе диагностический прогресс
часто важнее преждевременной победы: он позволяет отличить настоящий предел
механизма от случайного неуспеха конкретной настройки.

\begin{takeaway}
Главная задача этого свода -- не объявить задним числом, что серия ``доказала'',
а аккуратно восстановить логику движения. В ряде мест более ценно не то, что
некоторый метод оказался лучшим, а то, что удалось честно показать: текущий тип
механизма уперся в свой естественный потолок, и дальнейшее подкручивание тех же
ручек уже не обещает качественного скачка.
\end{takeaway}

\section*{Что именно считается экспериментом в этой книге}

Здесь под ``экспериментом'' понимается не только исполненный ноутбук. В конце
серии появляются и такие ветви, которые на момент свода существуют в разном
статусе:

\begin{itemize}[leftmargin=2em]
\item полностью исполненный ноутбук с сохраненными output-ячейками и отчетом;
\item исполненный ноутбук, но без окончательного markdown-репорта;
\item проектная развилка, подробно спроектированная как следующий опыт, но еще
не прогнанная.
\end{itemize}

Это особенно важно для \texttt{exp8b}. Его нельзя молча подать как уже
завершенный опыт. Но и выбрасывать его из исследовательской линии было бы
ошибкой: он играет роль концептуального фильтра, который должен был отделить
дефекты транспортной архитектуры от дефектов самой embedding-геометрии.

\section*{Как читать этот текст}

Есть два разумных маршрута чтения.

\begin{enumerate}[leftmargin=2em]
\item \textbf{Линейный маршрут.} Читать подряд, как историю исследования. Тогда
лучше начинать с введения и не перепрыгивать через главы, потому что поздние
серии (\texttt{7x--9a}) содержательно опираются на различия, выработанные в
\texttt{6x}.
\item \textbf{Маршрут по вопросам.} Если интересует одна конкретная тема --
например, PageRank-поправка, capacity, overlap cue-fields или semantic
proximity edges, -- можно идти сразу в соответствующую серию. Но в этом случае
нужно помнить, что поздние главы используют терминологию, введенную в первых
двух главах.
\end{enumerate}

\section*{Стилистическое замечание}

Язык текста намеренно педагогичен, но не упрощен до уровня популяризации.
Читатель здесь предполагается математически и алгоритмически подготовленным.
Поэтому базовые понятия вроде графа, случайного блуждания, косинусной
схожести или стационарного распределения не разжевываются с нуля; вместо этого
основной акцент сделан на том, как именно эти объекты начинают работать внутри
архитектуры памяти и почему в одном месте они помогают, а в другом создают
иллюзию прогресса.

\chapter{Исследовательская программа: граф, сигнал, метрики}

\section{Почему серия вообще начинается с эмбриона}

В архитектуре RIN графовая память не возникает как чистая абстракция. Она
должна рано или поздно стать рабочим субстратом, по которому проходит сигнал,
куда коммитятся новые следы и из которого извлекается контекст. Поэтому уже
начиная со второго эксперимента задача формулируется двойственно:

\begin{enumerate}[leftmargin=2em]
\item нужен \textbf{эмбриональный граф}, достаточно реальный, чтобы на нем
вылезали содержательные дефекты;
\item нужен \textbf{контролируемый стенд}, достаточно малый и читаемый, чтобы на
нем можно было измерять не только факт прохождения сигнала, но и качество
retrieval.
\end{enumerate}

Отсюда вся дальнейшая архитектура серии:

\begin{itemize}[leftmargin=2em]
\item реальные источники -- для столкновения с настоящей структурной грязью;
\item золотые синтетические миры -- для локализации конкретных механизмов;
\item перенос с золотого мира на эмбрион -- как проверка того, переживает ли
формула выход из лабораторной колбы.
\end{itemize}

\section{Как выглядит базовый графовый объект}

Во всех сериях используется одна и та же базовая идея памяти как направленного
взвешенного графа. Узлы представляют концепты, события, состояния, а позднее --
эпизодические или служебные элементы. Ребра несут тип и вес.

Минимально полезные поля узла:

\begin{itemize}[leftmargin=2em]
\item текстовое содержимое;
\item тип узла (\texttt{concept}, \texttt{episode}, \texttt{claim},
\texttt{relationship\_state} и т.\,д.);
\item embedding в общем векторном пространстве;
\item метаданные происхождения и эпистемического статуса.
\end{itemize}

Минимально полезные поля ребра:

\begin{itemize}[leftmargin=2em]
\item источник и приемник;
\item тип отношения (\texttt{IsA}, \texttt{RelatedTo}, \texttt{UsedFor},
\texttt{Causes}, \texttt{HasProperty} и пр.);
\item вес проводимости.
\end{itemize}

Уже на этом уровне скрыта важная методологическая ловушка. Один и тот же граф
можно рассматривать по меньшей мере в трех разных ролях:

\begin{enumerate}[leftmargin=2em]
\item как \textbf{хранилище фактов и эпизодов};
\item как \textbf{среду распространения сигнала};
\item как \textbf{объект переписывания и консолидации}.
\end{enumerate}

Серия \texttt{exp2--exp9a} почти целиком посвящена второй роли, но время от
времени касается первой и третьей. Из-за этого полезно заранее различать
операторы чтения и записи.

\section{Как выглядит retrieval в базовой версии}

Пусть задан набор стартовых якорей. Им присваивается начальная амплитуда, после
чего по графу запускается затухающее распространение. В простейшей форме можно
думать о шаге диффузии как о правиле

\[
a_{t+1}(v)=\lambda \sum_{u \to v} a_t(u)\, P(u \to v),
\]

где:

\begin{itemize}[leftmargin=2em]
\item $a_t(v)$ -- уровень активации узла $v$ на шаге $t$;
\item $\lambda$ -- коэффициент затухания;
\item $P(u \to v)$ -- нормированная доля потока, идущая из $u$ в $v$.
\end{itemize}

В реальных ноутбуках используется не буквальная марковская цепь с бесконечным
пределом, а конечный процесс с отсечением слабых хвостов. Практически это
значит следующее:

\begin{enumerate}[leftmargin=2em]
\item в стартовые узлы впрыскивается фиксированный заряд;
\item на каждом шаге активны только узлы выше минимального порога;
\item выходящий поток делится между исходящими ребрами пропорционально их весу;
\item процесс идет ограниченное число шагов или до опустошения фронта.
\end{enumerate}

Из этого уже видно, почему серия так быстро вышла к проблеме хабов. Узел может
оказаться ярким по двум разным причинам:

\begin{enumerate}[leftmargin=2em]
\item он действительно хорошо объясняется данным запросом;
\item в него удобно сходится много путей, даже если по смыслу он слишком общий.
\end{enumerate}

\section{Три уровня retrieval-скоринга}

По ходу серии постепенно кристаллизуется удобное разложение retrieval на три
содержательно разных уровня.

\subsection{Локальный транспорт}

Это вопрос о том, \emph{дошел ли сигнал вообще} и как он растекся по имеющимся
путям. Здесь живут:

\begin{itemize}[leftmargin=2em]
\item исходные якоря;
\item затухание;
\item веса ребер;
\item типизация проводимости;
\item возможные дополнительные каналы вроде semantic edges.
\end{itemize}

\subsection{Локальное удержание}

Даже если два узла получили одинаковый приток, не обязано быть так, что они
одинаково \emph{удерживают} этот приток как локально значимую активацию. Отсюда
рождаются знаменатели семейства $R(B)$:

\[
\retained(B)=\frac{\incoming(B)}{R(B)}.
\]

На этом уровне обсуждаются:

\begin{itemize}[leftmargin=2em]
\item степень узла;
\item взвешенная исходящая сила;
\item активный fan-in;
\item локальная capacity как ``эффективный объем представления''.
\end{itemize}

\subsection{Глобальная структурная слава}

Наконец, даже хороший локальный отклик может оказаться слишком высоким просто
потому, что узел глобально знаменит. Тогда появляется еще одна поправка:

\[
\score(v\mid q)=\frac{\text{LocalResponse}(v\mid q)}{\text{Prior}(v)^\eta}.
\]

Именно здесь возникают in-strength, out-strength, PageRank и их сравнение.

\begin{takeaway}
Одно из важнейших достижений серии состоит в том, что эти три уровня перестали
смешиваться. В начале работ любой неуспех мог быть списан на ``плохие
эмбеддинги'' или ``неудачную формулу''. К середине серии стало ясно, что
транспорт, удержание и глобальная слава -- разные механизмы, и провал любого из
них не лечится простой правкой соседнего.
\end{takeaway}

\section{Роль embedding-пространства}

Embedding в этой серии играет несколько ролей сразу.

\begin{enumerate}[leftmargin=2em]
\item Он нужен для anchoring: чтобы query-слова можно было связать с узлами
графа.
\item Он нужен для семантического gate: чтобы весить релевантность узла текущему
запросу.
\item Он нужен для additive bypass в \texttt{8a}: similarity начинает напрямую
входить в итоговый score.
\item Он нужен для semantic-edge construction в \texttt{9a}: ближайшие соседи в
embedding-пространстве становятся кандидатами на дополнительные ребра.
\end{enumerate}

Это делает особенно опасной ситуацию, когда embedding-пространство само по себе
дефектно. Если косинус живет в узком конусе, generic nodes могут становиться
псевдохабами еще до того, как сигнал вообще вошел в графовую динамику. Поэтому
ветка \texttt{8b} возникает не как побочная эстетика, а как попытка убрать
важный конфаундер: не следует ли приписываемый графу дефект на самом деле
свойствам самой similarity-геометрии.

\section{Какими метриками пользовалась серия}

Ранняя часть серии опиралась на простые и очень понятные критерии:

\begin{itemize}[leftmargin=2em]
\item правильный узел должен быть выше отвлекающего;
\item broad hub не должен побеждать специфический ответ, если запрос не о нем;
\item наоборот, legitimate hub не должен быть задавлен там, где запрос именно о
нем.
\end{itemize}

Позднее появляются более тонкие метрики:

\begin{itemize}[leftmargin=2em]
\item число пройденных pairwise inequality;
\item число решенных запросов;
\item margin между лучшим правильным и лучшим distractor;
\item mean rank и MRR;
\item manifold-aware success вместо проверки одного канонического узла;
\item специальные danger-checks: hub reflood и RAG collapse.
\end{itemize}

Сюда стоит добавить два методологических принципа, выработанных по ходу серии.

\begin{methodbox}
\textbf{Принцип 1.} Нельзя ранжировать методы по ``суммарному margin'', если они
по-разному масштабируют одну и ту же базовую картину. Иначе получится ложная
победа просто более агрессивной нормировки.

\textbf{Принцип 2.} Как только граф начинает выдавать не один канонический узел,
а соседний семантический manifold, оценка должна это учитывать. Иначе мы будем
объявлять провалом ситуацию, в которой память по существу попала в правильный
район, но не в ту формулировку вершины, которую человек заранее выписал в
таблицу ответов.
\end{methodbox}

\section{Почему понадобились золотые миры}

Реальный граф нужен для реализма, но он почти бесполезен как чистый измеритель
retrieval-качества, если там нет ground truth. Поэтому серия несколько раз
возвращается к hand-authored worlds. Это не ``игра в песочнице'' в плохом
смысле, а нормальный исследовательский ход:

\begin{enumerate}[leftmargin=2em]
\item на реальном графе вылавливаются инфраструктурные и содержательные дефекты;
\item на золотом мире локализуется механизм;
\item затем гипотеза переносится обратно на более реальный срез.
\end{enumerate}

Слабость золотого мира известна заранее: он зависит от авторского конструирования.
Его сила тоже известна заранее: только он позволяет быстро различить, является
ли неуспех следствием плохой формулы, плохого графа или грязной постановки.

\section{Как устроен дальнейший рассказ}

Оставшиеся главы идут не просто по номерам, а по логическим сериям.

\begin{itemize}[leftmargin=2em]
\item Главы \texttt{exp2--exp5} строят базовые объекты: эмбриональный граф,
золотой мир, семантический gate и write-path.
\item Глава о \texttt{6x} разбирает локальное удержание и глобальные priors.
\item Глава о \texttt{7x} показывает, как эти идеи переносятся на реальный
эмбрион и где именно перенос начинает ломаться.
\item Глава о \texttt{8a--9a} фиксирует смещение от ``оптимизации знаменателя'' к
проблеме транспортной архитектуры и реконструкции путей.
\item Заключительная глава собирает общую карту того, что серия действительно
доказала и где ее естественная развилка дальше.
\end{itemize}

\chapter{От эмбриона к золотому миру: \texttt{exp2} и \texttt{exp3}}

\section{\texttt{exp2}: эмбриональный граф из реальных источников}

\subsection{Смысл эксперимента}

\texttt{exp2} был первым настоящим столкновением с вопросом: можно ли вообще
собрать из реальных источников такой граф, который будет не просто лежать на
диске, а поддерживать полный Stage-1 pipeline. Речь шла не о красоте retrieval,
а о минимальной жизнеспособности всей линии:

\begin{enumerate}[leftmargin=2em]
\item seeded loading подграфа из реальных данных;
\item anchoring пользовательского запроса в этом графе;
\item диффузия активации;
\item упаковка результата в ограниченный контекст.
\end{enumerate}

Важное отличие от поздних экспериментов состоит в том, что здесь еще нет
полноценной ground-truth оценки качества ответа. Эксперимент отвечает на более
ранний вопрос: работает ли вообще ``физика системы'' на настоящем материале.

\subsection{Стенд}

Источники были следующими:

\begin{itemize}[leftmargin=2em]
\item полный дамп ConceptNet;
\item агрегированный ATOMIC 2019;
\item английский NumberBatch.
\end{itemize}

Загрузка строилась как seeded BFS от ограниченного набора входов с радиусом 3--4
и потолком числа concept-узлов. Дополнительно использовались:

\begin{itemize}[leftmargin=2em]
\item отдельный более жесткий порог для \texttt{RelatedTo};
\item точное лемматическое связывание ATOMIC с отключенным fuzzy fallback;
\item общий векторный словарь NumberBatch для якорения и similarity.
\end{itemize}

В терминах будущей архитектуры этот опыт проверял, может ли эмбрион
\emph{существовать как вычислимый объект}, а не лишь как схема в manual.

\subsection{Что здесь могло получиться}

На старте было две альтернативы.

\begin{enumerate}[leftmargin=2em]
\item \textbf{Оптимистическая.} Пайплайн в целом работает, а проблемы будут
касаться главным образом качества retrieval.
\item \textbf{Пессимистическая.} Реальные данные выявят системные дефекты еще до
момента, когда вообще можно говорить о retrieval-качестве: ошибки разметки,
схлопывание подписей, фантомные хабы, неработающие embedding lookup и т.\,п.
\end{enumerate}

\subsection{Что реально произошло}

Произошел более сильный и более полезный вариант пессимистической развилки:
pipeline в целом заработал, но именно на нем всплыл целый класс ошибок, которые
синтетический мир не мог бы показать.

\subsubsection*{(1) Ошибка labeling из POS-tagged URI}

Часть узлов получала содержимое вроде одной буквы \texttt{n}, потому что
последний сегмент URI принимался за label узла, хотя это был POS-tag. Ошибка
оказалась не косметической: в одном из запусков она затронула около трети узлов.

\subsubsection*{(2) Дублирование ребер}

Двунаправленное индексирование adjacency делало одно и то же ребро доступным из
обоих концов, а BFS потом повторно добавлял его в граф. Нужна была дедупликация
по assertion id самого дампа.

\subsubsection*{(3) Производство хабов связью \texttt{RelatedTo}}

Это одна из ключевых находок всей ранней серии. Оказалось, что проблема не в
``больших степенях вообще'', а в том, что именно \texttt{RelatedTo} на реальных
данных образует слишком широкие и семантически рыхлые маршруты. Отсюда уже в
\texttt{exp2} родилась мысль, что тип отношения в графе не является декоративной
меткой: он меняет реальную физику retrieval.

\subsubsection*{(4) Несовпадение формата NumberBatch и ConceptNet-URI}

English-only NumberBatch оказался записан как bare terms, а не как полные
ConceptNet URI. Поэтому все lookup молча проваливались, и embeddings не
подгружались вовсе.

\subsubsection*{(5) Переизбыточное ATOMIC-linking}

Даже после исправления NumberBatch слишком мягкая схема embedding-based linking
присоединяла к графу лавину эпизодов. Переход к exact lemma match дал резкое
сжатие числа episode nodes и тем самым показал, что на этом уровне важно не
``чуть лучше подобрать threshold'', а сначала вообще вырваться из режима
неконтролируемого фан-аута.

\subsection{Инфраструктурный результат}

После исправлений pipeline заработал end-to-end на реальных данных. Запрос
``what is the difference between lisbon and bacalhau'' успешно якорился,
распространял активацию и упаковывался в контекст.

Но этот результат сам по себе не означал, что retrieval хорош. Он означал более
скромную и при этом чрезвычайно важную вещь: архитектура вышла из стадии чистой
спекуляции и впервые столкнулась с реальной семантической грязью.

\begin{takeaway}
\texttt{exp2} доказал не качество памяти, а жизнеспособность эмбриона как
вычислимого объекта. Еще важнее он доказал, что реальные данные порождают
собственный класс багов, которых не существует на игрушечных графах. Это и
обосновывает всю дальнейшую линию: без реального эмбриона часть архитектурных
дефектов просто оставалась бы невидимой.
\end{takeaway}

\subsection{Что эксперимент оставил открытым}

Главный открытый вопрос после \texttt{exp2}: что значит ``хороший'' retrieval на
реальном графе, если там нет ground truth. Сомнительный ranking можно заметить,
но нельзя строго назвать неверным без контролируемой постановки.

Именно отсюда рождается \texttt{exp3}. Нужен был не новый retrieval law, а
маленький золотой мир, где правильный ответ заранее известен.

\section{\texttt{exp3}: золотой микромир как измеритель retrieval}

\subsection{Зачем понадобился второй стенд}

После \texttt{exp2} стало понятно, что два вопроса нельзя больше смешивать:

\begin{enumerate}[leftmargin=2em]
\item может ли pipeline работать на реальных данных;
\item выдает ли он правильный ответ, когда правильность вообще можно проверить.
\end{enumerate}

\texttt{exp3} не вводит новый retrieval-механизм. Он вводит новую
\emph{измерительную среду}: маленький hand-authored граф, где каждая неудача
может быть названа по имени.

\subsection{Конструкция золотого мира}

Граф был намеренно очень маленьким и семантически прозрачным. В нем были два
сюжетных кластера.

\begin{itemize}[leftmargin=2em]
\item Кофейный: \texttt{coffee}, \texttt{caffeine}, \texttt{beverage},
\texttt{hot drink}, \texttt{tea}, \texttt{juice}, \texttt{soda}.
\item Географический: \texttt{north dakota}, \texttt{south dakota},
\texttt{north}.
\end{itemize}

Смысл кластера про кофе был особенно важен. Он позволял сравнить:

\begin{itemize}[leftmargin=2em]
\item специфический, содержательно правильный ответ \texttt{caffeine};
\item широкий, структурно связный, но концептуально грубый distractor
\texttt{beverage}.
\end{itemize}

\subsection{Какие именно сценарии проверялись}

В \texttt{exp3} было четыре сценария.

\begin{enumerate}[leftmargin=2em]
\item \textbf{Multi-word anchoring.} Проверка, что \texttt{north dakota}
якорится как единый термин, а не распадается на \texttt{north} и
\texttt{dakota}.
\item \textbf{Distractor vs answer.} Ключевой сценарий:
``why coffee gives energy'' не произносит \texttt{caffeine} явно, поэтому
retrieval обязан выйти к нему за счет графовой структуры и similarity, а не за
счет банального подглядывания.
\item \textbf{Multi-anchor query.} Проверка, что разные темы могут жить вместе в
одном контексте.
\item \textbf{Regression on irrelevant words.} Проверка, что stopwords и
служебные слова не ломают anchoring.
\end{enumerate}

\subsection{Главный результат}

Сценарии 1, 3 и 4 прошли хорошо. Но второй сценарий дал первый по-настоящему
диагностический удар по retrieval-механике.

Для запроса про кофе и энергию peak activation выглядела так:

\begin{center}
\begin{tabular}{lcc}
\toprule
Узел & Амплитуда & Роль \\
\midrule
\texttt{beverage} & 146.67 & distractor \\
\texttt{coffee} & 100.00 & anchor \\
\texttt{hot drink} & 100.00 & anchor \\
\texttt{tea} & 100.00 & anchor \\
\texttt{caffeine} & 26.67 & правильный ответ \\
\bottomrule
\end{tabular}
\end{center}

Это не ``слегка неудачный ranking'', а качественно неправильное поведение:
широкий хаб побеждает специфический ответ более чем впятеро.

\subsection{Что именно это доказало}

Важен не только сам факт ошибки, но и ее структура. \texttt{caffeine} проиграл
не потому, что он был плохо представлен как embedding. Он проиграл потому, что
в базовой диффузии нет механизма, который бы делал различие между:

\begin{itemize}[leftmargin=2em]
\item узлом, в который сходится много путей просто по топологии;
\item узлом, который является более точным ответом на вопрос.
\end{itemize}

Это сразу породило рабочую гипотезу для \texttt{exp4}: возможно, хватает просто
добавить query-conditioned semantic gate, то есть умножать поток на релевантность
узла запросу.

\begin{riskbox}
Именно здесь серия могла бы свернуть в ложную сторону и начать считать, что
проблема целиком в embeddings. Но \texttt{exp3} устроен так, что дефект уже виден
на очень маленьком и прозрачном графе. Поэтому рабочая гипотеза после него --
не ``надо срочно менять векторное пространство'', а ``в operator diffusion
отсутствует релевантностный множитель''.
\end{riskbox}

\subsection{Связь \texttt{exp2} и \texttt{exp3}}

Пара \texttt{exp2}--\texttt{exp3} образует очень важное методологическое кольцо.

\begin{itemize}[leftmargin=2em]
\item \texttt{exp2} показал, что без реального контента не увидеть целый класс
инфраструктурных дефектов.
\item \texttt{exp3} показал, что без золотого мира нельзя всерьез говорить о
качестве retrieval, даже если pipeline уже технически работает.
\end{itemize}

Вместе они вводят основной стиль всей дальнейшей серии: чередование между
реализмом и контролируемостью.

\chapter{От семантического gate к write-path: \texttt{exp4} и \texttt{exp5}}

\section{\texttt{exp4}: semantic gate как первая попытка исправить retrieval}

\subsection{Что именно проверялось}

После \texttt{exp3} гипотеза выглядела почти естественно: если диффузия никак не
учитывает релевантность узла запросу, то, может быть, достаточно встроить в нее
множитель

\[
\mathrm{rel}(B\mid q),
\]

зависящий от similarity между узлом $B$ и query-векторами.

Тогда переходный вклад вдоль ребра будет определяться не только весом ребра, но
и тем, насколько приемник вообще подходит к смыслу запроса.

\subsection{Постановка}

Стенд сохранялся максимально близким к \texttt{exp3}:

\begin{itemize}[leftmargin=2em]
\item тот же hand-authored кофейный граф;
\item тот же запрос ``why coffee gives energy'';
\item те же embeddings;
\item единственное изменение -- наличие или отсутствие semantic gate.
\end{itemize}

Это был хороший чистый опыт: если ranking исправится, значит, проблема
действительно была только в том, что диффузия игнорировала query relevance.

\subsection{Что ожидалось при двух конкурирующих картинах}

\begin{enumerate}[leftmargin=2em]
\item \textbf{Картина A: причина в отсутствии relevance-gate.}
Тогда после добавления gate \texttt{caffeine} должен обогнать
\texttt{beverage}.
\item \textbf{Картина B: причина не только в relevance-gate, но и в самой
топологии сходимости путей.} Тогда gate улучшит картину, но не перевернет ее.
\end{enumerate}

\subsection{Результат}

Вторая картина оказалась верной. После введения gate распределение изменилось в
нужную сторону, но не настолько, чтобы исправить ranking:

\begin{center}
\begin{tabular}{lcc}
\toprule
Узел & Без gate & С gate \\
\midrule
\texttt{beverage} & 146.67 & 129.46 \\
\texttt{caffeine} & 26.67 & 27.20 \\
\texttt{juice} & 23.47 & 11.49 \\
\texttt{soda} & 23.47 & 13.65 \\
\bottomrule
\end{tabular}
\end{center}

То есть gate работал \emph{качественно правильно}, но \emph{количественно
недостаточно}.

\subsection{Почему это особенно важный результат}

Здесь произошла одна из самых тонких и самых полезных развязок всей ранней
серии. Можно было бы поспешно сказать:

\begin{quote}
``Ну вот, similarity плохая. Посмотрим на слово \texttt{energy} и увидим, что
\texttt{beverage} к нему даже ближе, чем \texttt{caffeine}. Значит, виновато
embedding-пространство''.
\end{quote}

Но аккуратный разбор показал, что это ложный след. По фактическому правилу
агрегации $\mathrm{rel}(B\mid q)$ узел \texttt{caffeine} \emph{и так} был
релевантнее \texttt{beverage}. Проигрыш происходил не на уровне per-node
relevance, а на уровне \emph{суммарной амплитуды после сходимости нескольких
путей}.

\begin{takeaway}
\texttt{exp4} показал очень важную вещь: можно иметь узел, который \emph{сам по
себе} более релевантен запросу, и все равно проигрывать, если конкурирующий узел
получает поток по нескольким сходящимся маршрутам. Это первый явный выход серии
к проблеме convergence topology.
\end{takeaway}

\subsection{Какой следующий вопрос возник немедленно}

После \texttt{exp4} стало недостаточно обсуждать только similarity. Требовался
эксперимент, который бы изолировал следующее утверждение:

\begin{quote}
``Если менять только топологию сходимости, оставляя запрос, embeddings и gate
неизменными, то ranking действительно должен переворачиваться''.
\end{quote}

Это и становится задачей \texttt{exp6}; но до него в серии появляется еще одна
ветвь -- \texttt{exp5}, посвященная не retrieval, а write-path.

\section{\texttt{exp5}: write-path как отдельный архитектурный срез}

\subsection{Почему эта глава здесь вообще стоит}

На первый взгляд \texttt{exp5} выбивается из серии. Вокруг него нет ни
\texttt{caffeine}, ни \texttt{PageRank}, ни topology correction. Но именно это и
важно: память не может состоять только из retrieval. Нужен механизм записи новых
узлов, дедупликации, коммита и форензического поиска по времени.

\texttt{exp5} поэтому играет роль боковой, но принципиально важной ветви: он
проверяет, что граф не только можно читать, но и можно \emph{корректно
модифицировать}.

\subsection{Постановка}

Сценарий был полностью контролируемым. Вместо автоматического извлечения
кандидатов из естественного текста вручную задавались объекты \texttt{Candidate},
имитирующие результат будущего экстрактора.

Сюжет -- ``поездка на Азоры'':

\begin{enumerate}[leftmargin=2em]
\item покупка плавок для поездки;
\item повтор того же события, который должен быть опознан как дубликат;
\item покупка зонта для той же поездки;
\item отдельное событие про ноутбук в другой день.
\end{enumerate}

Проверялись:

\begin{itemize}[leftmargin=2em]
\item корректный коммит узлов;
\item пропуск эпистемического дубликата;
\item forensic index по timestamp;
\item ситуация, где один кандидат ссылается на другой кандидат из того же batch.
\end{itemize}

\subsection{Какая альтернатива тут была на самом деле}

Хотя сценарий выглядит техническим, в нем скрыта содержательная архитектурная
развилка.

\begin{enumerate}[leftmargin=2em]
\item \textbf{Наивный single-pass commit.} Кандидаты коммитятся и сразу
разрешают свои ребра.
\item \textbf{Batch-aware two-pass commit.} Сначала весь пакет кандидатов
проходит gate и получает id, а уже потом разрешаются их внутренние ссылки.
\end{enumerate}

Если второй режим не нужен, значит внутренняя структура event-batch не играет
роли. Если нужен, значит write-path обязан уважать пакетность события.

\subsection{Что обнаружилось}

На первом же прогоне всплыл настоящий баг: если первый кандидат ссылается на
сущность, которую второй кандидат того же batch должен определить как полноценный
узел, single-pass логика создает stub и тем самым портит последующий commit.

В конкретном сценарии это означало, что \texttt{azores trip} оказывался не
\texttt{episode}-узлом, а generic \texttt{claim}-stub. Более того, правильный
кандидат затем ошибочно считался дубликатом уже существующего узла.

\subsection{Исправление}

Исправление состояло в переходе к двухпроходной логике:

\begin{enumerate}[leftmargin=2em]
\item весь batch проходит gating;
\item все одобренные кандидаты создают голые узлы;
\item внутри batch строится \texttt{content -> id} map;
\item только после этого разрешаются ребра, причем сначала по этой карте, и лишь
затем через механизм stub.
\end{enumerate}

\subsection{Результат после исправления}

После введения двухпроходного коммита сценарий прошел чисто:

\begin{itemize}[leftmargin=2em]
\item оба нужных узла коммитились корректно;
\item дубликат пропускался;
\item \texttt{azores trip} получал правильный \texttt{node\_kind};
\item forensic lookup по дню 9 находил только релевантные узлы и не подмешивал
ноутбук из другого дня.
\end{itemize}

\subsection{Почему этот опыт важен для всей книги}

\texttt{exp5} показывает, что у архитектуры памяти есть как минимум две разные
линии зрелости:

\begin{itemize}[leftmargin=2em]
\item retrieval-линия, где исследуется прохождение и ранжирование сигнала;
\item write-path-линия, где исследуется корректность коммита и структуры
обновления графа.
\end{itemize}

Именно благодаря этому \texttt{exp5} важно не потерять в общем своде. Без него
может возникнуть ложное впечатление, что серия занималась только чтением памяти.
На деле уже здесь проясняется, что даже простой commit требует не меньшей
аккуратности, чем retrieval.

\begin{openbox}
\texttt{exp5} оставил открытым следующий вопрос: что будет, когда ссылки начнут
идти не только внутри одного batch, но и между несколькими событиями, особенно
если ранний event сначала породил stub, а поздний должен его доопределить. Это
уже не проблема локального коммита, а проблема долговременной консолидации и
ревизии графа.
\end{openbox}

\section{Промежуточный итог после \texttt{exp4} и \texttt{exp5}}

К этому моменту исследовательская линия уже расщепилась на две зрелые оси.

\begin{enumerate}[leftmargin=2em]
\item Retrieval-ось показала, что простой semantic gate не решает проблему
слишком ярких широких узлов.
\item Write-path-ось показала, что даже controlled commit требует
batch-aware-процедуры и не выдерживает наивного single-pass режима.
\end{enumerate}

Дальнейшая серия \texttt{6x} возвращается именно к retrieval-оси и начинает
разбирать ее уже более математически: что именно удерживает сигнал, и почему
некоторые узлы светятся слишком ярко не по смыслу, а по топологии и глобальной
славе.

\chapter{Серия \texttt{6x}: от topology correction к гибриду local + global}

\section{Зачем вообще понадобилась серия \texttt{6x}}

После \texttt{exp4} стало ясно, что семантический gate сам по себе не вытягивает
специфический ответ из-под широкого хаба. Но это еще не говорило, \emph{какой
именно} механизм отсутствует. Серия \texttt{6x} строится как последовательное
расщепление вопроса:

\begin{enumerate}[leftmargin=2em]
\item действительно ли дело в топологии сходимости путей;
\item если да, можно ли исправить это локальной поправкой на удержание;
\item если локальная поправка помогает, какой именно класс знаменателей
осмыслен;
\item если локального удержания мало, нужен ли еще и глобальный structural
prior.
\end{enumerate}

Важно, что здесь впервые появляется уже почти ``физическая'' интерпретация
узла: не всякий inflow одинаково превращается в salience, потому что разные узлы
по-разному удерживают и рассеивают приходящий поток.

\section{\texttt{exp6}: topology-aware retrieval correction}

\subsection{Цель опыта}

\texttt{exp6} должен был сделать то, что \texttt{exp4} еще не мог: отделить
собственно semantic relevance от эффекта многократной сходимости на хабе.

Для этого опыт делился на две части:

\begin{enumerate}[leftmargin=2em]
\item topology surgery: меняется только структура путей;
\item correction ablation: исходный граф сохраняется, меняется только формула
поправки.
\end{enumerate}

\subsection{Часть A: хирургия топологии}

В первой части проверялась почти причинная гипотеза:

\begin{quote}
если проблема действительно в количестве и силе сходящихся маршрутов, то
достаточно удалить пару входов в \texttt{beverage} или добавить пару входов в
\texttt{caffeine}, и ranking должен перевернуться даже без смены embeddings и
query.
\end{quote}

Именно это и произошло.

\begin{center}
\begin{tabular}{lccc}
\toprule
Вариант графа & \texttt{caffeine} & \texttt{beverage} & Победа \\
\midrule
Исходный & 27.20 & 129.46 & нет \\
Срезать часть входов в \texttt{beverage} & 37.57 & 27.38 & да \\
Добавить входы в \texttt{caffeine} & 95.10 & 77.17 & да \\
Сделать оба изменения & 141.30 & 19.26 & да \\
\bottomrule
\end{tabular}
\end{center}

Это один из самых чистых результатов всей серии: при неизменном query и
неизменном semantic gate ranking переворачивается одной лишь правкой topology.

\subsection{Часть B: пост-аккумуляционные поправки}

После этого можно было уже предметно тестировать корректирующие формулы. На
исходном графе сравнивались:

\begin{itemize}[leftmargin=2em]
\item отсутствие поправки;
\item деление на $\deg(B)^\gamma$;
\item деление на $\fanin_t(B)^\gamma$.
\end{itemize}

Ключевым было то, что штраф накладывался \emph{после слияния вкладов} в узле
назначения. Это принципиально: если наказывать поток локально у источника, то
узел, у которого есть единственный путь наружу, все равно отправит весь свой
заряд в этот путь.

\subsection{Результат}

Единственной формулой, действительно перевернувшей ranking, оказалось деление
на полную степень узла:

\[
\incoming(B)\mapsto \frac{\incoming(B)}{\deg(B)}.
\]

Это был грубый, но убедительный proof of concept.

\subsection{Что эксперимент реально доказал}

\texttt{exp6} не доказал, что $1/\deg(B)$ есть закон памяти. Он доказал более
узкое, но очень важное утверждение:

\begin{quote}
если узел побеждает за счет чрезмерной сходимости путей, то поправка должна
действовать на merged accumulation, а не только на локальный выбор ребра у
источника.
\end{quote}

\section{\texttt{exp6b}: richer golden world и первое сравнение семейств}

\subsection{Почему одного кофейного fixture уже было мало}

Успех на одной конфигурации \texttt{caffeine} против \texttt{beverage} еще не
говорит, что найдена хорошая retrieval law. Можно легко подобрать формулу,
которая работает только на одном сценарии.

Поэтому \texttt{exp6b} строит богатый, но все еще читаемый золотой мир на 33
узлах. Его смысл не в ``реализме мира'', а в разнообразии структурных ситуаций.

\subsection{Четыре кластера графа}

В графе были выделены:

\begin{itemize}[leftmargin=2em]
\item домашний кластер (\texttt{apartment}, \texttt{wallet}, \texttt{keys},
\texttt{cats}, \texttt{feeding alarm}, \texttt{elevator turnaround});
\item travel cluster (\texttt{trip}, \texttt{airport}, \texttt{passport},
\texttt{boarding pass}, \texttt{hotel});
\item work cluster (\texttt{office}, \texttt{meeting}, \texttt{laptop},
\texttt{charger}, \texttt{dead laptop});
\item coffee/health cluster (\texttt{coffee}, \texttt{caffeine},
\texttt{beverage}, \texttt{energy}, \texttt{water}).
\end{itemize}

Принципиально важно, что запросы включали не только anti-hub cases, но и
legitimate-hub cases. Иначе любую агрессивную нормировку можно было бы ошибочно
объявить хорошей.

\subsection{Формулы}

Появляются семейства

\[
R_0(B)=1,
\]
\[
R_1(B)=1+\beta\, \degout(B)^\gamma,
\]
\[
R_2(B)=1+\beta\, \sout(B)^\gamma,
\]
\[
R_3(B)=\alpha+\beta\, \sout(B)^\gamma,
\]
\[
R_4(B)=\capacity(B)+\beta\, \sout(B)^\gamma,
\]
\[
R_5(B)=\alpha+\capacity(B)+\beta\, \sout(B)^\gamma.
\]

Здесь впервые появляется capacity $C(B)$ как локальная ширина представления.

\subsection{Что можно было ожидать}

Если ``утечка'' исчерпывает всю проблему, должен был победить leak-only класс
$R_2$ или $R_3$. Если же разные узлы отличаются не только по суммарной силе
исходящих каналов, но и по собственной концептуальной ширине, тогда capacity
должна дать добавочное преимущество.

\subsection{Результат}

В \texttt{exp6b} лучшим in-sample и leave-one-out победителем стало семейство

\[
R_4(B)=\capacity(B)+0.1\, \sout(B)^{0.5}.
\]

Оно решило все 7 запросов и дало наибольший общий margin. Базовый raw case
проваливал только 6/7, а leak-only формулы не смогли стать убедительными
победителями.

\subsection{Интерпретация}

Это был первый сильный аргумент в пользу следующего тезиса:

\begin{quote}
узел может быть широким в двух разных смыслах: потому что у него много исходящих
каналов и потому что он сам по себе представляет более объемное понятие.
\end{quote}

Именно эта двойственность и пыталась поймать пара ``capacity + leak''.

\begin{riskbox}
\texttt{exp6b} остается синтетическим опытом. Capacity в нем ручная. Поэтому
после него нельзя было честно говорить ``мы нашли закон памяти''. Можно было
говорить только ``семейство без capacity уже не выглядит достаточным''.
\end{riskbox}

\section{\texttt{exp6c}: mechanism benchmark для знаменателей}

\subsection{Зачем понадобился еще один шаг}

У \texttt{exp6b} оставалось методологическое возражение: возможно, ответ
приходил не столько из-за хорошего знаменателя, сколько из-за удачной geometry
anchors и semantic gate. Поэтому \texttt{exp6c} радикально очищает постановку:

\begin{itemize}[leftmargin=2em]
\item для каждого запроса заранее задан source event;
\item после этого диффузия идет без query-conditioned gating;
\item benchmark измеряет именно поведение знаменателя.
\end{itemize}

\subsection{Диагностические мотивы}

Здесь уже не большой смешанный мир, а набор маленьких мотивов, каждый из которых
стрессует отдельное различие:

\begin{itemize}[leftmargin=2em]
\item weather motif -- проверка weighted leakage;
\item kit motif -- наказание слишком грубого degree counting;
\item travel motif -- чистый capacity-only case;
\item health motif -- перенос исходного \texttt{caffeine}/\texttt{beverage}
конфликта в новый стенд.
\end{itemize}

\subsection{Результат}

В \texttt{exp6c} baseline и degree-control выступили слабо:

\begin{itemize}[leftmargin=2em]
\item $R_0$ решил 2/5 задач;
\item $R_1$ в leave-one-out вообще рухнул до 1/5;
\item weighted leak alone ($R_2$, $R_3$) улучшил ситуацию, но не дотянул до
полной стабильности;
\item capacity-aware families ($R_4$, $R_5$) оказались единственными по-настоящему
сильными кандидатами.
\end{itemize}

Лучшим leave-one-out вариантом стал

\[
R_5(B)=1.0+\capacity(B)+0.05\, \sout(B)^{0.5}.
\]

\subsection{Что добавил \texttt{exp6c} относительно \texttt{exp6b}}

\texttt{exp6b} говорил: ``capacity выглядит полезной''. \texttt{exp6c} добавил:
``capacity делает свою работу даже тогда, когда anchor geometry убрана из игры''.

Это уже не просто победа в смешанном benchmark, а настоящий mechanism audit.

\section{\texttt{exp6d}: глобальные priors поверх локального удержания}

\subsection{Почему после \texttt{6c} вопрос еще не был закрыт}

Даже если узел хорошо удерживает локальный приток, он может быть ярким просто
потому, что глобально структурно знаменит. Поэтому возникает следующий шаг:

\[
\score(v\mid q)=\frac{\text{LocalResponse}(v\mid q)}{\text{Prior}(v)^\eta}.
\]

Содержательный вопрос уже не в том, ``заменяет ли PageRank локальную физику'',
а в том, добавляет ли глобальный prior \emph{нечто не сводимое} к локальному
удержанию.

\subsection{Методы}

Сравнивались:

\begin{itemize}[leftmargin=2em]
\item $P0$: raw diffusion;
\item $P1$: raw / out-strength$^\eta$;
\item $P2$: raw / in-strength$^\eta$;
\item $P3$: raw / PageRank$^\eta$;
\item $P4$: local-only retained diffusion;
\item $P5$: local retained diffusion / PageRank$^\eta$.
\end{itemize}

При этом локальная часть в $P4$ и $P5$ фиксировалась на победителе из
\texttt{6c}.

\subsection{Ключевой мотив: ложная глобальная слава}

Особенно важным был morning motif, где \texttt{water} делали глобально известным
за счет множества дальних входящих связей, но локальный утренний cue делали почти
симметричным между \texttt{water} и \texttt{espresso}. Этот мотив проверял,
может ли prior отличить:

\begin{itemize}[leftmargin=2em]
\item ``узел яркий, потому что глобально везде встречается'';
\item ``узел яркий, потому что именно здесь он нужен''.
\end{itemize}

\subsection{Результат}

Лучшим методом стал гибрид

\[
P5(v\mid q)=\frac{\text{LocalResponse}(v\mid q)}{\PR(v)^{0.25}},
\]

где LocalResponse задается capacity-aware local denominator из \texttt{6c}.

Этот гибрид решил все 7 задач и in-sample, и leave-one-out.

\subsection{Что особенно важно в этом результате}

Успех \texttt{P5} означает не ``PageRank победил'', а нечто более тонкое:

\begin{enumerate}[leftmargin=2em]
\item purely global penalties полезны;
\item purely local retention полезна;
\item но наилучшее поведение дает их комбинация;
\item при этом $\eta$ должен быть мягким: полная PMI-подобная сила оказывается
слишком агрессивной.
\end{enumerate}

\begin{takeaway}
К концу серии \texttt{6x} складывается первая действительно зрелая retrieval
картина:
\[
\score(v\mid q)\approx
\frac{\incoming(v\mid q)}{R_{\mathrm{local}}(v)\,\PR(v)^\eta}.
\]
Локальная часть отвечает за удержание и противодействие smearing, глобальная --
за снятие structural fame. Это еще не перенос на реальный граф, но это уже
чистая и достаточно сильная рабочая гипотеза.
\end{takeaway}

\section{Что серия \texttt{6x} дала и чего не дала}

\subsection{Что она действительно установила}

\begin{enumerate}[leftmargin=2em]
\item Пост-аккумуляционная topology-aware correction имеет смысл.
\item Простая степень узла может работать как грубый control, но она слишком
тупа, чтобы считаться итоговым законом.
\item Weighted leak полезен, но сам по себе не исчерпывает задачу.
\item Capacity выполняет не декоративную, а содержательную работу.
\item Глобальный prior, особенно мягкий PageRank, добавляет нечто реальное поверх
локальной физики.
\end{enumerate}

\subsection{Чего серия еще не установила}

\begin{enumerate}[leftmargin=2em]
\item Что manual capacity переносима на реальный граф.
\item Что найденная формула не является артефактом hand-authored motifs.
\item Что hybrid local + global сохранит преимущество на шумном эмбриональном
подграфе.
\end{enumerate}

Именно поэтому после \texttt{6d} естественным следующим ходом становится не
очередной synthetic benchmark, а transfer-check на реальном seeded subgraph.

\chapter{Серия \texttt{7x}: перенос на эмбрион и локализация нового bottleneck}

\section{Почему перенос не был формальностью}

После \texttt{6d} легко было бы вообразить, что осталось просто применить
победивший гибрид к большому реальному графу. Но именно здесь и начинается самая
поучительная часть всей серии. Перенос показал, что формула, убедительная на
золотом мире, еще не гарантирует содержательного прогресса на эмбрионе. Причина
при этом оказалась не тривиальной.

\section{\texttt{exp7a}: end-to-end transfer check}

\subsection{Замысел}

\texttt{exp7a} спрашивал: сохраняются ли выводы \texttt{6c--6d} при переходе на
крупный seeded subgraph, где нет ручной capacity и где retrieval должен работать
на реальной, а не hand-authored топологии.

Сравнивались raw diffusion, in-strength prior, PageRank prior, local leak only и
hybrid local leak + PageRank.

\subsection{Стенд}

Подграф строился как seeded expansion по понятиям, продолжающим мотивы
\texttt{6x}: кофе, паспорт, apartment, charger, travel, morning и т.\,д.
Формально ceiling на concept nodes составлял 10\,000, но включение ATOMIC
создало гораздо более массивную оболочку episode nodes.

\subsection{Что получилось вместо чистого transfer-benchmark}

Фактический размер графа оказался таким:

\begin{itemize}[leftmargin=2em]
\item 10\,000 concept nodes;
\item 455\,481 episode nodes;
\item 465\,481 узел суммарно.
\end{itemize}

Уже одно это сильно меняет характер задачи. Но еще важнее был другой эффект:
automatic anchoring порождал шумные source clouds вроде \texttt{thing},
\texttt{matter}, \texttt{before}, \texttt{early}, \texttt{imbibement},
\texttt{morningness}.

\subsection{Результат}

Все tested methods дали одинаковую дискретную картину:

\begin{itemize}[leftmargin=2em]
\item 4/13 pairwise checks;
\item 3/7 solved queries.
\end{itemize}

Численные margins различались сильно, но ordering не менялся.

\subsection{Почему это не был ``нулевой'' опыт}

На первый взгляд такой tie можно было бы счесть провалом. Но на деле
\texttt{exp7a} дал один из важнейших диагнозов серии:

\begin{quote}
в реальном end-to-end режиме retrieval law почти не тестируется, потому что
pipeline уже ломается на этапе query $\to$ source injection.
\end{quote}

Кроме того, top-k были засорены echo самих source nodes, так что наблюдаемое
``retrieval'' часто было просто возвратом к уже впрыснутым якорям.

\begin{takeaway}
\texttt{exp7a} не опроверг local+global hybrid. Он показал, что на реальном
эмбрионе более ранний слой -- anchoring и source injection -- доминирует над
всем, что происходит позже. Это и стало основанием для controlled-source опыта
\texttt{7b}.
\end{takeaway}

\section{\texttt{exp7b}: controlled-source transfer}

\subsection{Почему это новый эксперимент, а не ``починка 7a''}

\texttt{7b} делает стандартный физический ход: замораживает вход, чтобы
проверять среду. Manual anchors здесь не подаются как финальная архитектура для
natural-language retrieval. Они используются как калиброванный впрыск сигнала.

Именно поэтому \texttt{7a} и \texttt{7b} должны читаться вместе:

\begin{itemize}[leftmargin=2em]
\item \texttt{7a} локализует bottleneck automatic anchoring;
\item \texttt{7b} временно убирает этот bottleneck из игры.
\end{itemize}

\subsection{Два варианта графа}

Отдельно сравнивались:

\begin{itemize}[leftmargin=2em]
\item ConceptNet only;
\item ConceptNet + ATOMIC.
\end{itemize}

Это было методически важно: после \texttt{7a} нельзя было молча допустить, что
гигантская ATOMIC shell никак не меняет topology priors.

\subsection{\texttt{7b.1}: первая controlled-source версия}

В первой версии manual cues были чистыми, но во многом еще scene-like:
\texttt{check in + airport}, \texttt{home + remember}, \texttt{morning +
breakfast}.

Результат: снова tie across methods.

\begin{itemize}[leftmargin=2em]
\item pairwise: 2/12;
\item solved: 2/7.
\end{itemize}

Главный урок этого захода: даже ручные anchors бесполезны, если сами cues
слишком общие и описывают сцену, а не механизм ответа.

\subsection{\texttt{7b.2}: diagnostic-cue version}

Во второй версии cues сделали более механизмо-ориентированными:

\begin{itemize}[leftmargin=2em]
\item \texttt{coffee + stimulant} вместо \texttt{coffee + energy};
\item \texttt{check in counter + document} вместо \texttt{check in + airport};
\item \texttt{wake up + tired} вместо \texttt{morning + breakfast}.
\end{itemize}

Это улучшило семантическую осмысленность top-k, но discrete separation все равно
не появилось.

\subsection{Результат \texttt{7b.2}}

На обеих graph variants получились одинаковые агрегаты:

\begin{itemize}[leftmargin=2em]
\item primary pairwise: 1/8;
\item primary solved: 1/5;
\item control pairwise: 2/4;
\item control solved: 1/2.
\end{itemize}

То есть controlled-source regime оказался методически правильным, но формулы
снова не разошлись.

\subsection{Что именно теперь стало ясно}

После \texttt{7b} уже нельзя было объяснять стагнацию одной лишь проблемой
automatic anchors. Bottleneck сместился глубже:

\begin{enumerate}[leftmargin=2em]
\item либо signal идет в broad intermediate hubs;
\item либо он попадает в graph-native adjacent manifold, а не в hand-named
target node;
\item либо оба эффекта происходят одновременно.
\end{enumerate}

Примеры этого были вполне содержательны:

\begin{itemize}[leftmargin=2em]
\item для morning query граф охотно выдавал \texttt{drink coffee}, а не обязательно
\texttt{coffee} или \texttt{espresso};
\item для laptop query граф выдавал \texttt{power}, \texttt{battery},
\texttt{generator}-family вместо \texttt{charger};
\item для check-in query граф поднимал procedural nodes вроде
\texttt{checking in} сильнее, чем \texttt{passport}.
\end{itemize}

\begin{riskbox}
Здесь очень легко сделать ложный вывод ``PageRank бесполезен'' или
``local leak не переносится''. На самом деле \texttt{7b} показал более узкую
вещь: оба механизма работают слишком поздно в pipeline, если transport already
lands in the wrong manifold.
\end{riskbox}

\section{\texttt{exp7c}: conjunctive cue-overlap}

\subsection{Новый вопрос}

После повторяющихся tie в \texttt{7b} естественно возникла новая гипотеза:
может быть, проблема не в знаменателе вообще, а в том, как складываются cue
fields. Если broad hub ярко светится от одного сильного cue, а правильный ответ
получает умеренную поддержку от нескольких независимых cues, то простой сумматор
может систематически переоценивать первый случай.

Отсюда и возникает overlap idea:

\[
\mathrm{Overlap}_{\mathrm{geom}}(v)=
\left(\prod_g (A_g(v)+\varepsilon)\right)^{1/G}.
\]

\subsection{Что проверялось}

Сравнивались:

\begin{itemize}[leftmargin=2em]
\item sum aggregation;
\item sum + PageRank;
\item geometric overlap;
\item geometric overlap + PageRank.
\end{itemize}

То есть это уже не еще один denominator benchmark, а новая попытка изменить
саму логику совместного учета нескольких cue-conditioned полей.

\subsection{Результат}

Именно здесь серия впервые довольно жестко исчерпала надежду на ``еще одну
формулу финального score''.

По дискретным метрикам sum-based и overlap-based варианты снова tied. В
исполненном \texttt{7c} это было зафиксировано как одинаковые primary outcomes
на обеих graph variants; в дальнейшей постановке \texttt{8a} этот итог
формулируется как устойчивый tie по уровням порядка \(4/10\) pairwise и \(3/6\)
solved для главных конкурентов.

\subsection{Почему это содержательно важно}

Overlap failure означает следующее:

\begin{quote}
нельзя ``перемножить'' путь, которого вообще нет. Если $A_g(v)=0$ для нужного
cue на нужной цели, то никакая интерференция или coincidence scoring не создаст
там содержательного сигнала.
\end{quote}

Это очень сильная отрицательная находка. Она переводит разговор от финального
скоринга к самой транспортной архитектуре.

\section{Что серия \texttt{7x} установила}

\subsection{Положительные результаты}

\begin{enumerate}[leftmargin=2em]
\item Automatic anchoring действительно является самостоятельным bottleneck.
\item Controlled-source regime нужен и методически оправдан.
\item Эффект source echo в retrieval-списках нужно явно устранять.
\item ATOMIC shell сильно влияет на размеры и priors, но не обязана менять сами
ordering outcomes.
\end{enumerate}

\subsection{Отрицательные результаты}

\begin{enumerate}[leftmargin=2em]
\item На реальном эмбрионе denominator families не дают такого separation, как на
золотых мирах.
\item Простая cue-overlap aggregation тоже не решает проблему.
\end{enumerate}

\subsection{Главный диагностический вывод}

К концу \texttt{7x} становится правдоподобной новая общая картина:

\begin{quote}
узкое место находится не столько в пост-фактум нормировке уже пришедшего
сигнала, сколько в том, что сигнал либо не доходит до правильного manifold,
либо доходит не теми путями, либо приходит в слишком соседний, а не точный
ответный узел.
\end{quote}

Отсюда уже естественным образом рождается серия \texttt{8a--9a}, где вопрос
смещается с denominators на transport wall.

\chapter{Серия \texttt{8a--9a}: транспортная стена, similarity и достройка путей}

\section{Почему после \texttt{7x} уже нельзя было просто искать еще один denominator}

К началу \texttt{8a} две надежды уже были достаточно честно проверены:

\begin{enumerate}[leftmargin=2em]
\item post-accumulation denominators сами по себе не отделяют методы на реальном
эмбрионе;
\item overlap cue-fields тоже не создает specificity там, где transport до цели
не доходит.
\end{enumerate}

Следовательно, следующий вопрос должен был звучать иначе:

\begin{quote}
может быть, проблема не в том, как пересчитать уже полученное поле, а в том, что
сигнал \emph{вообще идет не туда} или не доходит до нужного manifold?
\end{quote}

\section{\texttt{exp8a}: bypass и repair как две попытки пробить transport wall}

\subsection{Замысел}

\texttt{8a} сознательно отвернулся от еще одной ``кухни знаменателей'' и взялся
за два других рычага:

\begin{enumerate}[leftmargin=2em]
\item \textbf{bypass:} дать similarity прямой доступ в итоговый score, а не
держать ее только как gate внутри diffusion;
\item \textbf{repair:} ослабить смазывающие типы ребер, в первую очередь
\texttt{RelatedTo}, чтобы transport меньше стекал в catch-all zones.
\end{enumerate}

Появилась и еще одна методологическая новация: вместо проверки одного
канонического target node вводится \textbf{target manifold}. Это очень важный
сдвиг. С этого места серия перестает считать провалом любой случай, когда граф
выдал не exact label, а корректный соседний graph-native ответ.

\subsection{Методы}

Испытывались следующие варианты:

\[
M0(v)=\frac{transport(v)}{\PR(v)^\eta},
\]
\[
M_{\mathrm{gate}}(v)=transport_{\mathrm{gated}}(v),
\]
\[
M1(v)=\frac{transport(v)+\lambda\,\simf(v\mid q)}{\PR(v)^\eta},
\]
\[
M2(v)=\frac{transport_{\mathrm{routed}}(v)+\lambda\,\simf(v\mid q)}{\PR(v)^\eta},
\]
\[
M3(v)=\frac{M2(v)}{\text{genericness}(v)^\mu}.
\]

В качестве control сохранялся и overlap-вариант из \texttt{7c}.

\subsection{Почему сравнение было честным}

Особенно важно, что \(M0 \to M1\) меняет только одну вещь: additive similarity.
Это позволяет отделить старый semantic gate как ``legacy control'' от нового
вопроса: помогает ли similarity, если дать ей прямой канал в score.

\subsection{Что именно проверялось}

Primary suite была достаточно строгой:

\begin{itemize}[leftmargin=2em]
\item \texttt{coffee + stimulant} против \texttt{beverage}-family;
\item \texttt{check in + document} против procedural and generic distractors;
\item \texttt{laptop + battery} против \texttt{office}/\texttt{power};
\item \texttt{wake up + tired} против \texttt{morning}/\texttt{bed}.
\end{itemize}

Параллельно были control queries, где broad topic не должен был быть уничтожен
over-correction.

\subsection{Главный результат}

В своей окончательной, уже подправленной версии \texttt{8a} дал очень чистую
картину:

\begin{center}
\begin{tabular}{lcccc}
\toprule
Метод & primary manifold & strict & mean rank & control \\
\midrule
\texttt{M0} & 2/4 & 1/4 & 232.0 & 2/2 \\
\texttt{M\_gate} & 2/4 & 1/4 & 217.2 & 2/2 \\
\texttt{M1} & 2/4 & 1/4 & 63.5 & 2/2 \\
\texttt{M2} & 2/4 & 1/4 & 63.5 & 2/2 \\
\texttt{M3} & 2/4 & 1/4 & 135.8 & 2/2 \\
\texttt{CTRL\_overlap} & 1/4 & 0/4 & 5021.2 & 2/2 \\
\bottomrule
\end{tabular}
\end{center}

\subsection{Что из этого следует}

Здесь одновременно есть сильный отрицательный и умеренный положительный
результат.

\subsubsection*{Сильный отрицательный результат}

Overlap control теперь уже не просто ``не победил'', а стал явно хуже. Значит,
ветка ``многократное схождение cue-fields само по себе родит правильный
ответ'' может считаться практически исчерпанной на данном субстрате.

\subsubsection*{Умеренный положительный результат}

Additive similarity действительно улучшает rank structure. Это видно по падению
mean rank с 232.0 до 63.5 у \(M1\). Но она не увеличивает число дискретных
побед.

Иными словами, bypass способен подтягивать правильный manifold из глубины хвоста,
но пока не превращается в механизм надежного переворачивания hardest pairwise
cases.

\subsubsection*{Еще один отрицательный результат}

Routing через \texttt{RelatedTo} в текущем виде почти ничего не прибавил:
\(M2\) практически совпал с \(M1\). Genericness penalty в \(M3\) даже ухудшила
часть показателей.

\subsection{Качественное чтение отдельных запросов}

Именно top-k inspection делает выводы \texttt{8a} убедительными.

\begin{itemize}[leftmargin=2em]
\item В coffee-case \texttt{caffeine} действительно становится ближе, но над ним
остаются \texttt{served hot}, \texttt{coffee shop}, \texttt{made from coffee
beans}.
\item В passport-case система с additive similarity не находит \texttt{passport},
а проваливается в морфологические и procedural neighbours вроде
\texttt{documentor}, \texttt{photodocument}, \texttt{checking in}.
\item В charger-case signal по-прежнему вязнет в \texttt{power}-family и
соседнем structural noise.
\item Morning-case выглядит лучше других: \texttt{drink coffee} действительно
входит в сильный manifold, но не кристаллизуется как бесспорный победитель.
\end{itemize}

\begin{takeaway}
\texttt{8a} не нашел ``победную формулу'', но очень качественно расчистил поле.
После него уже можно уверенно говорить: проблема не решается ни overlap, ни еще
одним знаменателем; additive similarity полезна, но работает как мягкий rescue,
а не как полноценный retrieval engine.
\end{takeaway}

\section{\texttt{exp8b}: de-coning similarity как проектная развилка}

\subsection{Статус}

На момент этого свода \texttt{8b} существует как тщательно продуманная ветка, но
не как исполненный ноутбук. Это нужно проговорить прямо.

\subsection{Почему эта ветка возникла}

Как только similarity из \texttt{8a} становится first-class term, встает вопрос:
а хороша ли вообще сама similarity. Если embedding-space анизотропно и живет в
узком конусе, то cosine margins изначально сжаты, а generic nodes становятся
embedding-hubs.

Тогда часть дефекта, который мы приписываем transport, на самом деле может
оказаться следствием грязной геометрии \texttt{NumberBatch}.

\subsection{Что именно предполагалось проверить}

\texttt{8b} должен был сравнить:

\begin{itemize}[leftmargin=2em]
\item raw cosine;
\item mean-centering;
\item all-but-the-top;
\item CSLS;
\item per-query z-score.
\end{itemize}

Причем сначала на intrinsic критериях (pairwise cosine distribution, hubness,
margin widening), а затем как plug-in replacement для \(\simf(v\mid q)\) в
лучшем score из \texttt{8a}.

\subsection{Почему даже неисполненная ветка все равно важна}

Потому что она задает очень важное ограничение для \texttt{9a}: если
semantic-proximity edges строить на сыром косинусе в узком конусе, можно просто
переписать embedding-hubness в topology и тем самым ухудшить граф.

\begin{openbox}
Если бы \texttt{8b} дал сильный положительный результат, это сделало бы
de-coned similarity стандартным строительным материалом и для additive bypass, и
для semantic-edge construction. Если бы не дал, это сняло бы важный конфаундер и
позволило бы еще увереннее винить transport rather than cosine. Именно поэтому
ветка проектно значима, даже еще не будучи прогнанной.
\end{openbox}

\section{\texttt{exp9a}: dual-geometry transport и semantic-proximity edges}

\subsection{Большой вопрос}

Если \texttt{8a} показал, что additive similarity лишь подтягивает правильные
узлы, но не решает проблему содержательно, то возникает более радикальная идея:
может быть, в графе просто \emph{нет нужных путей}, и их надо аккуратно
достроить.

Отсюда и dual-geometry picture:

\begin{itemize}[leftmargin=2em]
\item graph geometry: typed relational edges;
\item embedding geometry: proximity in vector space.
\end{itemize}

Новая идея состояла в том, чтобы добавить minority channel semantic edges:
mutual-kNN на embeddings, но с mutual filter, degree caps и throttled
conductance.

\subsection{Постановка}

Испытывались семейства:

\begin{itemize}[leftmargin=2em]
\item \(G0b\): no edges, no bypass;
\item \(G1,G2,G3\): только semantic edges с растущим \(w_{\mathrm{sem}}\);
\item \(G0a\): realistic baseline из \texttt{8a} с additive similarity;
\item \(G1r,G2r,G3r\): edges + bypass.
\end{itemize}

Параллельно измерялись два danger checks:

\begin{itemize}[leftmargin=2em]
\item hub reflood: generic distractors не должны снова полезть наверх;
\item RAG collapse: ranking не должен выродиться в pure embedding-kNN.
\end{itemize}

\subsection{Главная сводка}

Результат по основным retrieval-запросам оказался таким:

\begin{center}
\begin{tabular}{lccc}
\toprule
Метод & \((w_{\mathrm{sem}}, \lambda)\) & primary pairwise & median rank \\
\midrule
\texttt{G0b} & (0.0, 0.0) & 1/3 & 249 \\
\texttt{G1} & (0.1, 0.0) & 1/3 & 127 \\
\texttt{G2} & (0.3, 0.0) & 1/3 & 90 \\
\texttt{G3} & (1.0, 0.0) & 1/3 & 43 \\
\texttt{G0a} & (0.0, 1.0) & 1/3 & 111 \\
\texttt{G1r} & (0.1, 1.0) & 1/3 & 98 \\
\texttt{G2r} & (0.3, 1.0) & 1/3 & 59 \\
\texttt{G3r} & (1.0, 1.0) & 0/3 & 42 \\
\bottomrule
\end{tabular}
\end{center}

\subsection{Как это читать}

Здесь снова есть важное разделение между rank-improvement и task-solution.

\begin{itemize}[leftmargin=2em]
\item Да, semantic edges заметно подтягивают target manifolds вверх по рангу.
\item Нет, они не увеличивают число дискретных побед над strongest distractor.
\item При слишком сильном semantic-channel combined regime даже портится
(см. \(G3r\)).
\end{itemize}

\subsection{Сильная сторона \texttt{9a}}

Особенно интересным оказался bridge probe на alias/synonym pairs:

\begin{itemize}[leftmargin=2em]
\item \texttt{cat -> kitty}: rank 1161 \(\to\) 285 \(\to\) 278;
\item \texttt{couch -> sofa}: 1488 \(\to\) 37 \(\to\) 22.
\end{itemize}

То есть там, где проблема действительно состоит в отсутствии короткого моста
между очень близкими представлениями, semantic channel делает реальную работу.

\subsection{Почему этого оказалось недостаточно}

В сложных retrieval cases, особенно passport-case, semantic edges начинали
соединять не ``правильный ответный смысл'', а лексико-морфологическое облако:
\texttt{checking in}, \texttt{check in counter}, \texttt{documentor},
\texttt{photodocument}, \texttt{redocument} и т.\,п.

Это очень сильный диагностический результат. Он показывает, что blind
embedding-proximity repair шьет graph gaps, но шьет их не по той семантике,
которая нужна ассоциативной памяти.

\subsection{Danger checks}

Здесь ноутбук тоже оказался очень полезен.

\begin{itemize}[leftmargin=2em]
\item Hub reflood не взорвался на умеренных \(w_{\mathrm{sem}}\): median rank
generic distractors даже немного рос.
\item Но RAG-collapse стал заметен: overlap top-20 с pure-sim ranking вырос
примерно от 0.10 до 0.52--0.53.
\end{itemize}

Это значит, что сильный semantic channel начинает подменять граф векторным
retrieval. Память при этом не столько улучшается, сколько тайком становится
другим механизмом.

\begin{takeaway}
\texttt{9a} не дал структурного ``спасения'' линии, но очень хорошо локализовал
предел blind semantic-edge repair. Он показал, что у эмбриона действительно не
хватает мостов, однако строить их как глобальные mutual-kNN links по сырой
NumberBatch-геометрии недостаточно: такие мосты тянут не только к смысловым
соседям, но и к морфологическим и лексическим паразитным облакам.
\end{takeaway}

\section{Что серия \texttt{8a--9a} установила}

\subsection{Отрицательные результаты}

\begin{enumerate}[leftmargin=2em]
\item Overlap cue-fields не является недостающим ингредиентом retrieval.
\item Ни relation-type routing в текущем виде, ни genericness penalty не дали
решающего эффекта.
\item Blind semantic-proximity edges не превращают эмбрион в хорошую память.
\end{enumerate}

\subsection{Положительные и диагностические результаты}

\begin{enumerate}[leftmargin=2em]
\item Additive similarity полезна как мягкий bypass: она улучшает rank structure
и поднимает цель из глубокого хвоста.
\item Semantic edges годятся как средство alias/synonym bridging.
\item Но в содержательных commonsense retrieval cases главная проблема остается
глубже: signal либо идет по неправильным typed routes, либо graph-native answer
organizes itself not as the named target node.
\end{enumerate}

\subsection{Почти окончательный вывод этой линии}

К концу \texttt{9a} становится правдоподобной уже достаточно смелая мысль:

\begin{quote}
в текущем виде диффузионная retrieval-парадигма, даже с priors, leak laws,
overlap and semantic repair, может быть уперлась в свой естественный потолок на
данном субстрате.
\end{quote}

Это не означает, что диффузия бесполезна. Это означает, что от нее, возможно,
требуют того, что ближе к path-selective reasoning или к архитектуре
консолидации графа, чем к простому spreading process.

\chapter{Синтез: что доказала серия \texttt{exp2--exp9a} и куда она вывела}

\section{Четыре слоя выводов}

Чтобы не смешивать разные типы результата, полезно собрать серию по слоям.

\subsection{Слой 1: инфраструктура и эмбрион}

Этот слой строят прежде всего \texttt{exp2} и \texttt{exp5}.

\begin{itemize}[leftmargin=2em]
\item Реальный эмбриональный граф можно собрать и прогонять через retrieval
pipeline.
\item Реальные источники порождают собственный класс ошибок: URI-labeling,
\texttt{RelatedTo}-hubness, embedding-lookup mismatch, ATOMIC fan-out.
\item Write-path требует batch-aware commit semantics; наивный single-pass
ломается даже на малом сценарии.
\end{itemize}

\subsection{Слой 2: локальная механика retrieval}

Этот слой задается \texttt{exp3}, \texttt{exp4}, \texttt{exp6}, \texttt{exp6b},
\texttt{exp6c}.

\begin{itemize}[leftmargin=2em]
\item Простой diffusion ranking действительно может предпочесть широкий хаб
специфическому ответу.
\item Semantic gate исправляет это лишь частично.
\item Причина не сводится к плохой similarity: topology of convergence сама по
себе умеет задавить более релевантный узел.
\item Post-accumulation correction имеет физический смысл.
\item Weighted leakage полезен, но не исчерпывает картину.
\item Capacity как ``эффективный объем представления'' выполняет реальную
работу на диагностических стендах.
\end{itemize}

\subsection{Слой 3: глобальная структурная слава}

Этот слой кристаллизуется в \texttt{exp6d}.

\begin{itemize}[leftmargin=2em]
\item Local response и global fame -- разные механизмы.
\item Мягкий PageRank prior действительно помогает, если применяется поверх
локально осмысленного response.
\item Полная PMI-подобная сила prior оказывается слишком агрессивной.
\end{itemize}

\subsection{Слой 4: перенос на реальный эмбрион}

Этот слой строят \texttt{7a--9a}.

\begin{itemize}[leftmargin=2em]
\item Automatic anchoring -- самостоятельный bottleneck.
\item Controlled-source regime нужен, но даже он не возвращает separation между
формулами автоматически.
\item Graph-native answer manifold на реальном субстрате часто не совпадает с
ручным canonical target.
\item Neither denominator tweaking, nor overlap, nor blind semantic-edge repair
дают решающего breakthrough.
\end{itemize}

\section{Что именно можно считать уже опровергнутым}

Серия достаточно уверенно отсекла несколько популярных, но ложных надежд.

\begin{enumerate}[leftmargin=2em]
\item \textbf{``Надо просто добавить semantic gate''} -- нет, этого мало.
\item \textbf{``Проблема только в embeddings''} -- нет, topology of convergence
делает самостоятельную работу.
\item \textbf{``Достаточно подобрать хороший denominator''} -- на реальном
эмбрионе это не разделяет методы так, как на золотых мирах.
\item \textbf{``Нужно просто поощрять overlap cue-fields''} -- нет, overlap не
создает путь там, где его нет.
\item \textbf{``Можно достроить все мосты blind kNN edges и получить решение''}
-- нет, тогда память начинает подтягивать морфологическое облако и опасно
съезжает к vector-RAG.
\end{enumerate}

\section{Что можно считать реально поддержанным}

Здесь нужна аккуратность. Поддержанность не означает окончательную истину, но
означает, что серия дала сильный положительный или диагностически устойчивый
сигнал в пользу гипотезы.

\begin{enumerate}[leftmargin=2em]
\item \textbf{Typed graph matters.} Тип ребра меняет физику retrieval;
\texttt{RelatedTo} нельзя считать нейтральным проводником.
\item \textbf{Post-accumulation viewpoint is correct.} Если дефект заключается в
слиянии множества вкладов на broad hub, штрафовать нужно именно merged
destination response.
\item \textbf{Local and global corrections are orthogonal.} Local retention и
PageRank-like prior решают разные задачи.
\item \textbf{Source injection must be modeled explicitly.} На реальном графе
anchoring не является технической мелочью; он архитектурно равноправен diffusion
и scoring.
\item \textbf{Target manifold is the right late-stage evaluation object.} Когда
графовая память возвращает соседний, но осмысленный узел, ошибочно считать это
тем же типом провала, что и победу пустого хаба.
\end{enumerate}

\section{Где сейчас настоящий bottleneck}

Если собрать все вместе, то bottleneck выглядит уже не как ``плохая формула''.
Скорее он состоит из трех связанных уровней.

\subsection{(1) Качество графового субстрата}

Эмбрион сейчас слишком смешан:

\begin{itemize}[leftmargin=2em]
\item lexical artifacts;
\item morphologically noisy neighborhoods;
\item неоднородные типы отношений;
\item неполные мосты между реально близкими представлениями;
\item слишком рыхлая интеграция episode shell.
\end{itemize}

\subsection{(2) Качество source injection}

Даже controlled-source regime понадобился не случайно. Это значит, что архитектура
естественного query-to-anchor mapping еще не дошла до состояния, где последующие
эксперименты можно интерпретировать без оговорок.

\subsection{(3) Предел самого spreading-подхода}

Есть все больше оснований думать, что для части задач память хочет не просто
``куда растечется активация'', а что-то ближе к:

\begin{itemize}[leftmargin=2em]
\item path-sensitive retrieval;
\item typed route preference;
\item consolidation shortcuts;
\item maybe small proof-like or script-like selection поверх поля активации.
\end{itemize}

\section{Почему серия все же продвинулась далеко}

Иногда длинная череда отрицательных и диагностических результатов производит
впечатление застоя. Но в этой серии произошло обратное: исследовательский объект
стал значительно яснее.

В начале речь шла почти обо всем сразу: реальный граф, embeddings, retrieval,
write-path, gate, priors. К концу серии можно уже сказать заметно точнее:

\begin{itemize}[leftmargin=2em]
\item проблема не в том, что ``ничего не работает'';
\item проблема не сводится к одному плохо выбранному коэффициенту;
\item есть рабочие локальные механизмы;
\item есть рабочие диагностические стенды;
\item есть честно исчерпанные ветки;
\item и есть новая, уже более содержательная развилка для следующего цикла.
\end{itemize}

\section{Что логично делать дальше}

Из самой серии вытекают по меньшей мере три разумные программы продолжения.

\subsection{Программа A: мета-аудит и смена класса retrieval}

Если принять, что spreading-подход почти уперся в потолок, следующий шаг --
формализовать это и перейти к path-selective retrieval. Тогда новая серия должна
прямо спрашивать:

\begin{quote}
какая причинно-смысловая траектория лучше всего объясняет запрос, а не какой
узел просто ярче всех после диффузии.
\end{quote}

\subsection{Программа B: архитектура консолидации эмбриона}

\texttt{9a} очень сильно намекнул, что графу не хватает внутренних мостов, но
достраивать их blind kNN-способом нельзя. Значит, нужна не просто edge repair, а
полноценная программа консолидации:

\begin{itemize}[leftmargin=2em]
\item alias-links;
\item synonym bridges;
\item shortcut links по часто используемым путям;
\item controlled reverse-links;
\item co-activation based strengthening;
\item фильтрация морфологического мусора.
\end{itemize}

\subsection{Программа C: развести слои памяти}

Еще один естественный ход -- перестать ожидать от одного плоского графа, что он
одновременно будет:

\begin{itemize}[leftmargin=2em]
\item словарем;
\item commonsense memory;
\item episode memory;
\item ассоциативной сетью синонимов и вариантов;
\item рабочей средой диффузии.
\end{itemize}

Тогда следующая архитектура должна быть многослойной, а transport и retrieval --
межслойными.

\section{Финальная оценка серии}

Если сформулировать общий итог одной фразой, он звучит так:

\begin{quote}
серия \texttt{exp2--exp9a} не нашла окончательную retrieval law для RIN, но она
очень качественно расчленила проблему и показала, где заканчивается сила
нынешнего диффузионного подхода и откуда должен начинаться следующий архитектурный
шаг.
\end{quote}

Это уже большой результат. Более того, без него любой следующий шаг рисковал бы
снова раствориться в хаотическом подборе коэффициентов и частных эвристик.

\begin{takeaway}
Самое ценное, что дала серия, -- не отдельная формула, а новая ясность. Теперь
можно продолжать не с надеждой, что ``еще один знаменатель все спасет'', а с
довольно четким пониманием: где нам нужен другой механизм чтения, где -- другой
механизм роста графа, а где -- более зрелая организация самой памяти как
многослойной архитектуры.
\end{takeaway}


\end{document}
